\documentclass[12pt,a4paper]{article}
\usepackage[top=2.5cm,bottom=2.5cm,left=2.5cm,right=2.5cm,headheight=15pt]{geometry}
\usepackage{amsmath,amssymb}
\usepackage{booktabs,array,longtable,tabularx,multirow}
\usepackage{xcolor}
\usepackage{enumitem}
\usepackage{fancyhdr}
\usepackage{titlesec}
\usepackage{mdframed}
\usepackage{tcolorbox}
\tcbuselibrary{skins,breakable}
\usepackage{hyperref}

%% ---- Colors ----
\definecolor{planblue}{RGB}{0,82,165}
\definecolor{lightblue}{RGB}{220,235,255}
\definecolor{darkorange}{RGB}{180,70,0}
\definecolor{lightorange}{RGB}{255,235,210}
\definecolor{darkgreen}{RGB}{20,110,20}
\definecolor{lightgreen}{RGB}{200,240,205}
\definecolor{darkred}{RGB}{150,20,20}
\definecolor{lightred}{RGB}{255,210,210}
\definecolor{lightgray}{RGB}{242,242,242}
\definecolor{darkteal}{RGB}{0,105,92}
\definecolor{lightteal}{RGB}{200,230,225}

%% ---- tcolorbox styles ----
\newtcolorbox{qbox}[2][]{enhanced,breakable,
  colback=lightorange,colframe=darkorange,arc=4pt,
  title={\textbf{#2}},fonttitle=\bfseries\color{white},
  attach boxed title to top left={yshift=-2mm,xshift=4mm},
  boxed title style={colback=darkorange,arc=3pt},#1}

\newtcolorbox{solbox}[2][]{enhanced,breakable,
  colback=lightblue,colframe=planblue,arc=4pt,
  title={\textbf{#2}},fonttitle=\bfseries\color{white},
  attach boxed title to top left={yshift=-2mm,xshift=4mm},
  boxed title style={colback=planblue,arc=3pt},#1}

\newtcolorbox{keybox}[2][]{enhanced,breakable,
  colback=lightgreen,colframe=darkgreen,arc=4pt,
  title={\textbf{#2}},fonttitle=\bfseries\color{white},
  attach boxed title to top left={yshift=-2mm,xshift=4mm},
  boxed title style={colback=darkgreen,arc=3pt},#1}

\newtcolorbox{algobox}[2][]{enhanced,breakable,
  colback=lightgray,colframe=gray!60,arc=4pt,
  title={\textbf{#2}},fonttitle=\bfseries,
  attach boxed title to top left={yshift=-2mm,xshift=4mm},
  boxed title style={colback=gray!60,arc=3pt},#1}

\newtcolorbox{warnbox}[2][]{enhanced,breakable,
  colback=lightred,colframe=darkred,arc=4pt,
  title={\textbf{#2}},fonttitle=\bfseries\color{white},
  attach boxed title to top left={yshift=-2mm,xshift=4mm},
  boxed title style={colback=darkred,arc=3pt},#1}

%% ---- Page style ----
\pagestyle{fancy}\fancyhf{}
\lhead{\color{planblue}\textbf{Planning (STRIPS + POP) --- CSE 713 AI}}
\rhead{\color{planblue}Exam: 10 Jun 2026}
\cfoot{--- \thepage\ ---}
\renewcommand{\headrulewidth}{0.4pt}

%% ---- Section style ----
\titleformat{\section}{\large\bfseries\color{planblue}}{}{0em}{}[\titlerule]
\titleformat{\subsection}{\normalsize\bfseries\color{darkorange}}{}{0em}{}
\titleformat{\subsubsection}{\normalsize\bfseries\color{darkteal}}{}{0em}{}

%% ---- Document ----
\begin{document}

\begin{center}
{\LARGE\bfseries\color{planblue} CSE 713 --- Artificial Intelligence}\\[4pt]
{\large\bfseries STRIPS + Partial-Order Planning (Block World)}\\[2pt]
{\large Complete Past Paper Solutions: 2020--2024}\\[6pt]
{\normalsize Source: \textit{Planning.pdf} (50 slides) + \textit{plan2.pdf} (38 slides)}\\[2pt]
{\normalsize\color{darkred}\textbf{Every question from 2020--2024 included. Zero omissions.}}
\end{center}

\tableofcontents
\newpage

%% ==========================================
\section*{Topic Overview (from Slides)}
%% ==========================================

\begin{keybox}{STRIPS + POP --- What to Know Cold}
\textbf{STRIPS (Stanford Research Institute Problem Solver):}
\begin{itemize}[noitemsep]
  \item Action = Preconditions + Add-list + Delete-list
  \item State = conjunction of positive ground literals
  \item Closed World Assumption: anything not stated is false
  \item Solves the Frame Problem: only explicit effects change the state
\end{itemize}
\vspace{4pt}
\textbf{Partial-Order Planning (POP):}
\begin{itemize}[noitemsep]
  \item Plan = $(A, O, L)$: set of actions, ordering constraints, causal links
  \item Causal link: $A_p \xrightarrow{Q} A_c$ means $A_p$ achieves condition $Q$ for $A_c$
  \item Threat: action $A_t$ threatens link $(A_p, Q, A_c)$ if $A_t$ deletes $Q$ and could fall between $A_p$ and $A_c$
  \item Resolution: \textbf{Demotion} ($A_t \prec A_p$) or \textbf{Promotion} ($A_c \prec A_t$)
  \item Solution: no open conditions, no threats, consistent ordering
  \item POP is \textbf{sound and complete}
\end{itemize}
\vspace{4pt}
\textbf{Special actions for POP:}
\begin{itemize}[noitemsep]
  \item $A_0$: no preconditions; effects = initial state; must be \textbf{first}
  \item $A_\infty$: no effects; preconditions = goal; must be \textbf{last}
\end{itemize}
\end{keybox}

%% ==========================================
\section*{Block World STRIPS Operators (Memorise These)}
%% ==========================================

\begin{algobox}{4 Standard Block World Operators}
\begin{tabular}{@{}p{3.8cm}p{5.5cm}p{2.5cm}p{2.5cm}@{}}
\toprule
\textbf{Action} & \textbf{Preconditions} & \textbf{Add} & \textbf{Delete} \\
\midrule
PICKUP(b) & ONTABLE(b), CLEAR(b), ARMEMPTY & HOLDING(b) & ONTABLE(b), CLEAR(b), ARMEMPTY \\
PUTDOWN(b) & HOLDING(b) & ONTABLE(b), CLEAR(b), ARMEMPTY & HOLDING(b) \\
UNSTACK(b, c) & ON(b,c), CLEAR(b), ARMEMPTY & HOLDING(b), CLEAR(c) & ON(b,c), CLEAR(b), ARMEMPTY \\
STACK(b, c) & HOLDING(b), CLEAR(c) & ON(b,c), ARMEMPTY, CLEAR(b) & HOLDING(b), CLEAR(c) \\
\bottomrule
\end{tabular}

\medskip
\textit{Memory rule:} UNSTACK = pick up a block from another block.
PICKUP = pick up from table. STACK = place on a block. PUTDOWN = place on table.
\end{algobox}

\begin{warnbox}{The Standard 4-Block World (2022, 2023, 2024 --- appears verbatim every year)}
\textbf{Initial State:}
ON(B,A), ONTABLE(A), ONTABLE(C), ONTABLE(D), ARMEMPTY, CLEAR(B), CLEAR(C), CLEAR(D)\\
\emph{(A is NOT clear --- B is sitting on top of it.)}

\medskip
\textbf{Goal State:} ON(C,A), ON(B,D), ONTABLE(A), ONTABLE(D)

\medskip
\textbf{Total-order solution:}
UNSTACK(B,A) $\rightarrow$ STACK(B,D) $\rightarrow$ PICKUP(C) $\rightarrow$ STACK(C,A)
\end{warnbox}

\newpage

%% ==========================================
\section{2024 --- Q5 (9 marks)}
%% ==========================================

\subsection{Part (a) --- STRIPS + Household Robot Action Schema [4 marks]}

\begin{qbox}{2024 Q5(a) --- STRIPS Definition + Action Schema}
What is STRIPS? How are actions represented using STRIPS action schema?
Develop a planning problem using STRIPS defining one action for a
\textbf{Household Robot} (cleaning a room, turning on appliances, fetching objects)
with Preconditions, Add-effects, and Delete-effects.
\end{qbox}

\begin{solbox}{Solution}
\textbf{STRIPS (Stanford Research Institute Problem Solver)} is a planning framework
representing the world as a set of positive ground literals (propositions).
An action in STRIPS is an \textbf{action schema} consisting of three components:
\begin{enumerate}[noitemsep]
  \item \textbf{Preconditions}: a conjunction of positive, function-free literals that
    must hold in the current state before the action can execute.
  \item \textbf{Add-list (Add-effects)}: propositions that become \textbf{true} after the action.
  \item \textbf{Delete-list (Delete-effects)}: propositions that become \textbf{false} after the action.
\end{enumerate}

\textbf{Key property (Frame Problem solution):}
Only propositions explicitly in the add/delete lists change; everything else persists.

\bigskip
\textbf{Household Robot --- Three Action Schemas:}

\medskip
\texttt{Action: CLEAN\_ROOM(room)}\\
\phantom{xxx}Preconditions: $\{$IN(robot, room), DIRTY(room)$\}$\\
\phantom{xxx}Add-effects: $\{$CLEAN(room)$\}$\\
\phantom{xxx}Delete-effects: $\{$DIRTY(room)$\}$

\medskip
\texttt{Action: TURN\_ON(appliance, loc)}\\
\phantom{xxx}Preconditions: $\{$IN(robot, loc), AT(appliance, loc), OFF(appliance)$\}$\\
\phantom{xxx}Add-effects: $\{$ON(appliance)$\}$\\
\phantom{xxx}Delete-effects: $\{$OFF(appliance)$\}$

\medskip
\texttt{Action: FETCH(obj, loc)}\\
\phantom{xxx}Preconditions: $\{$IN(robot, loc), LOCATED(obj, loc)$\}$\\
\phantom{xxx}Add-effects: $\{$HOLDING(obj)$\}$\\
\phantom{xxx}Delete-effects: $\{$LOCATED(obj, loc)$\}$

\medskip
\emph{(Any one of the above three is sufficient for full marks --- the question asks for
``one action.'')}
\end{solbox}

\subsection{Part (b) --- POP for Standard Block World [5 marks]}

\begin{qbox}{2024 Q5(b) --- Partial-Order Planning}
Block World (Start / Goal as above).
Using Partial-Order Planning (POP):
\begin{enumerate}[noitemsep]
  \item Construct the minimal POP plan.
  \item List all causal links and ordering constraints.
  \item Explain how POP avoids backtracking compared to total-order planning.
\end{enumerate}
\end{qbox}

\begin{solbox}{Solution}
\textbf{Step 1: Actions required and their roles.}

\begin{itemize}[noitemsep]
  \item Goal \textbf{ON(C,A)} requires STACK(C,A). Its preconditions: HOLDING(C) and CLEAR(A).
  \item \textbf{CLEAR(A)} is not in the initial state (B is on A). Need UNSTACK(B,A) first.
  \item UNSTACK(B,A) produces HOLDING(B). Use it immediately: STACK(B,D) achieves goal ON(B,D).
  \item STACK(B,D) restores ARMEMPTY $\Rightarrow$ PICKUP(C) can run next.
  \item PICKUP(C) produces HOLDING(C) needed by STACK(C,A).
\end{itemize}

\textbf{Plan Steps:}
\[
S_1 = \text{UNSTACK}(B,A),\quad S_2 = \text{STACK}(B,D),\quad
S_3 = \text{PICKUP}(C),\quad S_4 = \text{STACK}(C,A)
\]

\bigskip
\textbf{Step 2: Causal Links.}

\begin{center}
\begin{tabular}{clc}
\toprule
\textbf{Producer} & \textbf{Condition} & \textbf{Consumer} \\
\midrule
$A_0$ & ON(B,A) & $S_1$ \\
$A_0$ & CLEAR(B) & $S_1$ \\
$A_0$ & ARMEMPTY & $S_1$ \\
$S_1$ & HOLDING(B) & $S_2$ \\
$A_0$ & CLEAR(D) & $S_2$ \\
$S_2$ & ARMEMPTY & $S_3$ \\
$A_0$ & ONTABLE(C) & $S_3$ \\
$A_0$ & CLEAR(C) & $S_3$ \\
$S_3$ & HOLDING(C) & $S_4$ \\
$S_1$ & CLEAR(A) & $S_4$ \\
$S_2$ & ON(B,D) & $A_\infty$ \\
$S_4$ & ON(C,A) & $A_\infty$ \\
$A_0$ & ONTABLE(A) & $A_\infty$ \\
$A_0$ & ONTABLE(D) & $A_\infty$ \\
\bottomrule
\end{tabular}
\end{center}

\textbf{Step 3: Ordering Constraints.}
\[
A_0 \prec S_1 \prec S_2 \prec S_3 \prec S_4 \prec A_\infty
\]
Justifications:
\begin{itemize}[noitemsep]
  \item $S_1 \prec S_2$: $S_2$ needs HOLDING(B) produced by $S_1$.
  \item $S_2 \prec S_3$: $S_3$ needs ARMEMPTY restored by $S_2$.
  \item $S_1 \prec S_4$: $S_4$ needs CLEAR(A) produced by $S_1$.
  \item $S_3 \prec S_4$: $S_4$ needs HOLDING(C) produced by $S_3$.
\end{itemize}

\textbf{Threat Analysis:} $S_1$ deletes ARMEMPTY, but $A_0 \xrightarrow{\text{ARMEMPTY}} S_1$
uses $S_1$ as its own consumer --- no external threat.
$S_3$ needs ARMEMPTY from $S_2$ (not $A_0$); $S_1 \prec S_2$ ensures $S_1$
cannot intervene between $S_2$ and $S_3$. \textbf{No unresolved threats.}

\bigskip
\textbf{Step 4: Why POP Avoids Backtracking.}
\begin{enumerate}[noitemsep]
  \item \textbf{Least Commitment Principle:} POP does not impose a total order upfront.
    It adds ordering constraints only when a causal link is threatened.
    Subgoals ON(C,A) and ON(B,D) are worked on independently.
  \item \textbf{No Goal Interaction Backtrack:} A total-order planner committing to
    ``achieve ON(C,A) first'' would stack C on A. Then achieving ON(B,D)
    requires unstacking B from A, which destroys the causal link for CLEAR(A).
    POP detects threats to causal links before execution and resolves them
    via promotion or demotion --- \emph{without restarting the search}.
  \item \textbf{Explicit Causal Support:} Every precondition has a named producer.
    If a new action threatens a link, POP either promotes the consumer or
    demotes the threat. This replaces backtracking with a targeted fix.
\end{enumerate}
\end{solbox}

\newpage

%% ==========================================
\section{2023 --- Q5 (9 marks)}
%% ==========================================

\subsection{Part (a) --- STRIPS + Air Cargo Action Schema [4 marks]}

\begin{qbox}{2023 Q5(a) --- STRIPS + Air Cargo}
What is STRIPS? What is the Action Schema of STRIPS?
Devise a planning problem for \textbf{Air Cargo Transport}
(loading cargo, unloading cargo, flying planes between airports).
\end{qbox}

\begin{solbox}{Solution}
\textbf{STRIPS:} See 2024 Q5(a) above.

\textbf{Air Cargo Transport --- Action Schemas:}

\medskip
\texttt{Action: LOAD(cargo, plane, airport)}\\
\phantom{xxx}Preconditions: $\{$AT(cargo, airport), AT(plane, airport)$\}$\\
\phantom{xxx}Add-effects: $\{$IN(cargo, plane)$\}$\\
\phantom{xxx}Delete-effects: $\{$AT(cargo, airport)$\}$

\medskip
\texttt{Action: UNLOAD(cargo, plane, airport)}\\
\phantom{xxx}Preconditions: $\{$IN(cargo, plane), AT(plane, airport)$\}$\\
\phantom{xxx}Add-effects: $\{$AT(cargo, airport)$\}$\\
\phantom{xxx}Delete-effects: $\{$IN(cargo, plane)$\}$

\medskip
\texttt{Action: FLY(plane, from, to)}\\
\phantom{xxx}Preconditions: $\{$AT(plane, from), AIRPORT(from), AIRPORT(to)$\}$\\
\phantom{xxx}Add-effects: $\{$AT(plane, to)$\}$\\
\phantom{xxx}Delete-effects: $\{$AT(plane, from)$\}$

\medskip
\textbf{Example Problem Instance:}\\
Initial: AT(C1, JFK), AT(P1, JFK), AIRPORT(JFK), AIRPORT(SFO)\\
Goal: AT(C1, SFO)\\
Plan: LOAD(C1, P1, JFK) $\rightarrow$ FLY(P1, JFK, SFO) $\rightarrow$ UNLOAD(C1, P1, SFO)
\end{solbox}

\subsection{Part (b) --- POP for Standard Block World [5 marks]}

\begin{qbox}{2023 Q5(b) --- POP Block World (same problem as 2024)}
Block World (same Start / Goal). Solve using partial-order planning.
\end{qbox}

\begin{solbox}{Solution}
The problem is identical to 2024 Q5(b). The minimal POP plan is:
\[
A_0 \prec S_1\!=\!\text{UNSTACK}(B,A) \prec S_2\!=\!\text{STACK}(B,D)
\prec S_3\!=\!\text{PICKUP}(C) \prec S_4\!=\!\text{STACK}(C,A) \prec A_\infty
\]

Causal links and ordering constraints are identical to 2024 Q5(b) --- see that section.

\textbf{POP vs.\ total-order planning:}
\begin{itemize}[noitemsep]
  \item \textbf{Total-order} commits to a sequence upfront. If goals interact
    (one goal's achievement undoes another's precondition), it backtracks
    and tries a different ordering --- potentially exponential search.
  \item \textbf{POP} defers ordering commitments. Goals ON(C,A) and ON(B,D)
    are treated as independent open conditions. The ordering
    $S_1 \prec S_4$ (UNSTACK before STACK(C,A)) is only imposed because
    $S_4$ needs CLEAR(A) produced by $S_1$ --- not because of an arbitrary
    total-order guess. No backtracking is needed.
\end{itemize}
\end{solbox}

\newpage

%% ==========================================
\section{2022 --- A-3 (9 marks)}
%% ==========================================

\subsection{Part (a) --- Block World Operator Preconditions [2 marks]}

\begin{qbox}{2022 A-3(a) --- Preconditions of Block World Operators}
Write the preconditions of:
UNSTACK(A,B),\ STACK(A,B),\ PICKUP(A),\ PUTDOWN(A).
\end{qbox}

\begin{solbox}{Solution --- Memorise These Verbatim}
\begin{itemize}[noitemsep,leftmargin=*]
  \item \textbf{UNSTACK(A, B):} $\{$ON(A,B),\ CLEAR(A),\ ARMEMPTY$\}$\\
    \hfill\emph{A is on top of B; A has nothing above it; arm is free.}
  \item \textbf{STACK(A, B):} $\{$HOLDING(A),\ CLEAR(B)$\}$\\
    \hfill\emph{Arm holds A; B is clear to receive A.}
  \item \textbf{PICKUP(A):} $\{$ONTABLE(A),\ CLEAR(A),\ ARMEMPTY$\}$\\
    \hfill\emph{A is on the table; A has nothing above it; arm is free.}
  \item \textbf{PUTDOWN(A):} $\{$HOLDING(A)$\}$\\
    \hfill\emph{Arm is holding A.}
\end{itemize}
\end{solbox}

\subsection{Part (b) --- Block World with STRIPS / POP [5 marks]}

\begin{qbox}{2022 A-3(b) --- Block World (same problem)}
Block World (same Start / Goal).
Solve using goal stack planning / partial-order planning / STRIPS.
\end{qbox}

\begin{solbox}{Solution --- Goal Stack / STRIPS Approach}
\textbf{Goal Stack Planning (STRIPS divide-and-conquer):}

\medskip
\begin{tabular}{@{}llp{7cm}@{}}
\toprule
\textbf{Step} & \textbf{Action} & \textbf{State after action} \\
\midrule
0 & --- (initial) &
  ON(B,A), ONTABLE(A), ONTABLE(C), ONTABLE(D),
  ARMEMPTY, CLEAR(B), CLEAR(C), CLEAR(D) \\
1 & UNSTACK(B,A) &
  ONTABLE(A), ONTABLE(C), ONTABLE(D),
  HOLDING(B), CLEAR(A), CLEAR(C), CLEAR(D) \\
2 & STACK(B,D) &
  ONTABLE(A), ONTABLE(C), ONTABLE(D), ON(B,D),
  ARMEMPTY, CLEAR(A), CLEAR(B), CLEAR(C) \\
3 & PICKUP(C) &
  ONTABLE(A), ONTABLE(D), ON(B,D),
  HOLDING(C), CLEAR(A), CLEAR(B) \\
4 & STACK(C,A) &
  ONTABLE(A), ONTABLE(D), ON(B,D), ON(C,A),
  ARMEMPTY, CLEAR(B), CLEAR(C) \\
\bottomrule
\end{tabular}

\medskip
\textbf{Goal verification:}
ON(C,A)~\checkmark\quad ON(B,D)~\checkmark\quad ONTABLE(A)~\checkmark\quad ONTABLE(D)~\checkmark

\medskip
\textbf{Why this order?}
UNSTACK(B,A) must come first because STACK(C,A) requires CLEAR(A),
which is blocked by B. After removing B, we immediately place it on D (achieving
ON(B,D)) and restore ARMEMPTY. Then PICKUP(C) $\rightarrow$ STACK(C,A) achieves ON(C,A).

\medskip
\emph{For the POP representation with causal links and ordering constraints, see 2024 Q5(b).}
\end{solbox}

\subsection{Part (c) --- Short Notes: Generative AI [2 marks]}

\begin{qbox}{2022 A-3(c) --- Generative AI}
Write short notes on Generative AI.
\end{qbox}

\begin{solbox}{Solution}
\textbf{Generative AI} refers to AI systems that create new content (text, images, code,
audio) statistically resembling their training data.

\begin{itemize}[noitemsep]
  \item \textbf{Key Models:} Large Language Models (LLMs like GPT, Claude) predict the
    next token; Diffusion models (DALL-E, Stable Diffusion) denoise from Gaussian noise;
    GANs pit a generator against a discriminator.
  \item \textbf{Applications:} Text/code generation, image synthesis, drug discovery,
    synthetic training data, automated planning.
  \item \textbf{Limitations:} Hallucinations (plausible but false outputs), training data
    bias, no persistent world-model or memory across sessions, high compute cost.
  \item \textbf{Relevance to AI planning:} LLMs can act as approximate forward planners,
    but lack the formal guarantees (soundness, completeness) of STRIPS/POP.
\end{itemize}
\end{solbox}

\newpage

%% ==========================================
\section{2021 --- Q6(b): Sussman Anomaly [4 marks]}
%% ==========================================

\begin{qbox}{2021 Q6(b) --- Sussman Anomaly with POP}
Block World --- \textbf{Sussman Anomaly}:\\[4pt]
\textbf{Initial State:} C on A, B on table, A on table\\
\phantom{xxxxxx}$\Rightarrow$ ON(C,A), ONTABLE(A), ONTABLE(B), CLEAR(C), CLEAR(B)\\[4pt]
\textbf{Goal State:} A on B, B on C\\
\phantom{xxxxxx}$\Rightarrow$ ON(A,B), ON(B,C)\\[4pt]
Develop an effective and complete plan using POP / STRIPS,
addressing threats to causal links and open conditions.
\end{qbox}

\begin{solbox}{Solution}
\textbf{Why it is the ``Sussman Anomaly'':}
No total-order planner can solve this without interleaving subplans.
\begin{itemize}[noitemsep]
  \item To achieve ON(A,B): put A on B. But B needs to be on C first, or A's position
    on C blocks things.
  \item To achieve ON(B,C): put B on C. But C currently has A on top of it.
  \item Goals cyclically interfere --- total-order planning backtracks endlessly.
    POP resolves this via threat detection.
\end{itemize}

\bigskip
\textbf{Block World Actions (simplified --- no arm required):}

\medskip
\textbf{A3: Move-C-from-A-to-Table}\\
\phantom{xxx}PRE: $\{$ON(C,A), CLEAR(C)$\}$ \quad ADD: $\{$ONTABLE(C), CLEAR(A)$\}$ \quad DEL: $\{$ON(C,A)$\}$

\medskip
\textbf{A1: Move-B-from-Table-to-C}\\
\phantom{xxx}PRE: $\{$ONTABLE(B), CLEAR(B), CLEAR(C)$\}$ \quad ADD: $\{$ON(B,C)$\}$ \quad DEL: $\{$ONTABLE(B), CLEAR(C)$\}$

\medskip
\textbf{A2: Move-A-from-Table-to-B}\\
\phantom{xxx}PRE: $\{$ONTABLE(A), CLEAR(A), CLEAR(B)$\}$ \quad ADD: $\{$ON(A,B)$\}$ \quad DEL: $\{$ONTABLE(A), CLEAR(B)$\}$

\bigskip
\textbf{POP Step-by-Step:}

\textbf{Step 1 --- Initial partial plan:} $A_0 \prec A_\infty$\\
$A_0$ effects: $\{$ON(C,A), ONTABLE(A), ONTABLE(B), CLEAR(B), CLEAR(C)$\}$\\
$A_\infty$ open conditions: $\{$ON(A,B), ON(B,C)$\}$

\medskip
\textbf{Step 2 --- Address ON(B,C):}\\
Add action A1. Causal link: $A_1 \xrightarrow{\text{ON(B,C)}} A_\infty$.\\
A1's open conditions: ONTABLE(B), CLEAR(B), CLEAR(C) --- all from $A_0$:\\
$A_0 \xrightarrow{\text{ONTABLE(B)}} A_1$,\quad
$A_0 \xrightarrow{\text{CLEAR(B)}} A_1$,\quad
$A_0 \xrightarrow{\text{CLEAR(C)}} A_1$.

\medskip
\textbf{Step 3 --- Address ON(A,B):}\\
Add action A2. Causal link: $A_2 \xrightarrow{\text{ON(A,B)}} A_\infty$.\\
A2's open conditions: ONTABLE(A), CLEAR(A), CLEAR(B).\\
$A_0 \xrightarrow{\text{ONTABLE(A)}} A_2$, \quad $A_0 \xrightarrow{\text{CLEAR(B)}} A_2$.\\
CLEAR(A) is \textbf{not} in $A_0$'s effects (C is on A) $\Rightarrow$ need action A3.

\medskip
\textbf{Step 4 --- Address CLEAR(A) for A2:}\\
Add action A3. Causal link: $A_3 \xrightarrow{\text{CLEAR(A)}} A_2$.\\
A3's open conditions: ON(C,A), CLEAR(C) --- from $A_0$:\\
$A_0 \xrightarrow{\text{ON(C,A)}} A_3$, \quad $A_0 \xrightarrow{\text{CLEAR(C)}} A_3$.

\medskip
\textbf{Step 5 --- Detect and Resolve Threats:}

\begin{itemize}[leftmargin=*]
  \item \textbf{Threat 1:} $A_2$ deletes CLEAR(B).\\
    Causal link $A_0 \xrightarrow{\text{CLEAR(B)}} A_1$ is threatened by $A_2$.\\
    $A_2$ could be placed between $A_0$ and $A_1$, deleting CLEAR(B) before $A_1$ uses it.\\
    \textbf{Resolution --- Promotion:} force $A_1 \prec A_2$
    (consumer $A_1$ before threat $A_2$). \checkmark

  \item \textbf{Threat 2:} $A_1$ deletes CLEAR(C).\\
    Causal link $A_0 \xrightarrow{\text{CLEAR(C)}} A_3$ is threatened by $A_1$.\\
    $A_1$ could be placed between $A_0$ and $A_3$, deleting CLEAR(C) before $A_3$ uses it.\\
    \textbf{Resolution --- Promotion:} force $A_3 \prec A_1$
    (consumer $A_3$ before threat $A_1$). \checkmark
\end{itemize}

\bigskip
\textbf{Final POP Plan (ordering):}
\[
A_0 \prec A_3 \prec A_1 \prec A_2 \prec A_\infty
\]

\textbf{Final action sequence:}
\begin{center}
\fbox{Move-C-from-A-to-Table} $\longrightarrow$
\fbox{Move-B-from-Table-to-C} $\longrightarrow$
\fbox{Move-A-from-Table-to-B}
\end{center}

\bigskip
\textbf{Complete Causal Link Table:}
\begin{center}
\begin{tabular}{clcl}
\toprule
\textbf{Producer} & \textbf{Condition} & \textbf{Consumer} & \textbf{Note} \\
\midrule
$A_0$ & ON(C,A) & $A_3$ & \\
$A_0$ & CLEAR(C) & $A_3$ & \emph{protected: $A_3 \prec A_1$} \\
$A_0$ & ONTABLE(B) & $A_1$ & \\
$A_0$ & CLEAR(B) & $A_1$ & \emph{protected: $A_1 \prec A_2$} \\
$A_0$ & CLEAR(C) & $A_1$ & \emph{(A0 also provides CLEAR(C) to A1)} \\
$A_3$ & CLEAR(A) & $A_2$ & \\
$A_0$ & ONTABLE(A) & $A_2$ & \\
$A_0$ & CLEAR(B) & $A_2$ & \\
$A_1$ & ON(B,C) & $A_\infty$ & \\
$A_2$ & ON(A,B) & $A_\infty$ & \\
\bottomrule
\end{tabular}
\end{center}
\end{solbox}

\newpage

%% ==========================================
\section{2020 --- B-Q6 (estimated 9 marks)}
%% ==========================================

\begin{qbox}{2020 B-Q6 --- STRIPS + Block World Planning}
\textbf{[Exact question not captured in question map.
Based on 2020--2024 pattern, this question covered STRIPS + Block World.]}
\end{qbox}

\begin{solbox}{Expected Content Based on Pattern}
The 2020 paper (Group-B Question 6) most likely asked:
\begin{enumerate}[noitemsep]
  \item \textbf{Define STRIPS} (preconditions, add-list, delete-list) + action schema.
  \item \textbf{Block World problem} --- either the Sussman Anomaly style
    (C on A, B on table $\rightarrow$ A on B, B on C) or the standard 4-block version.
  \item Possibly: forward/backward planning comparison or definition of planning as search.
\end{enumerate}

\textbf{Key preparation:} Be able to write both block world solutions cold:
\begin{itemize}[noitemsep]
  \item \textbf{4-block (2022--2024):}
    UNSTACK(B,A) $\to$ STACK(B,D) $\to$ PICKUP(C) $\to$ STACK(C,A)
  \item \textbf{Sussman Anomaly (2021):}
    Move-C-off-A $\to$ Move-B-onto-C $\to$ Move-A-onto-B
\end{itemize}
\end{solbox}

%% ==========================================
\section*{Quick Summary --- Exam Patterns}
%% ==========================================

\begin{keybox}{What to Memorise for 10 Jun Exam}
\textbf{ALWAYS asked:}
\begin{enumerate}[noitemsep]
  \item \textbf{STRIPS definition} (2--3 lines): preconditions + add-list + delete-list.
    Frame problem solved by: only explicit effects change state.
  \item \textbf{New domain action schema}: Household Robot (2024), Air Cargo (2023),
    generic (2022). Pattern: define 1--3 actions with PRE/ADD/DEL.
  \item \textbf{Block World solution} (same start/goal 2022--2024):\\
    UNSTACK(B,A) $\to$ STACK(B,D) $\to$ PICKUP(C) $\to$ STACK(C,A).
  \item \textbf{POP / causal links / ordering constraints} (2024, 2023, 2022): see table above.
  \item \textbf{Sussman Anomaly POP} (2021): $A_0 \to A_3 \to A_1 \to A_2 \to A_\infty$.
\end{enumerate}

\medskip
\textbf{Two threats in Sussman (memorise exactly):}
\begin{itemize}[noitemsep]
  \item $A_2$ deletes CLEAR(B) $\Rightarrow$ threatens $A_0\xrightarrow{\text{CLEAR(B)}}A_1$
    $\Rightarrow$ Promotion: $A_1 \prec A_2$.
  \item $A_1$ deletes CLEAR(C) $\Rightarrow$ threatens $A_0\xrightarrow{\text{CLEAR(C)}}A_3$
    $\Rightarrow$ Promotion: $A_3 \prec A_1$.
\end{itemize}

\medskip
\textbf{Operator preconditions (2022 asked verbatim):}
\begin{itemize}[noitemsep]
  \item UNSTACK(A,B): ON(A,B), CLEAR(A), ARMEMPTY
  \item STACK(A,B): HOLDING(A), CLEAR(B)
  \item PICKUP(A): ONTABLE(A), CLEAR(A), ARMEMPTY
  \item PUTDOWN(A): HOLDING(A)
\end{itemize}

\medskip
\textbf{POP vs total-order (asked 2024, 2023):}
POP uses \textit{least commitment} --- only orders actions when a causal link is threatened
(promotion/demotion). Total-order guesses a sequence and backtracks when goals interfere.
POP never needs to restart; it resolves threats locally.
\end{keybox}

\end{document}
